% chapters/04/section-05.tex
\section{Mengukur dan Menyetel Laju}
\label{sec:mengukur-laju}

Pernyataan pada Bagian~\ref{sec:mengendalikan-laju} bahwa laju tanpa
pembatasan memboroskan prosesor sebaiknya tidak diterima begitu saja.
Bagian ini memakai program pengukuran untuk membuktikannya sendiri, sekaligus
melatih kebiasaan yang penting sepanjang semester: mengukur, bukan menebak.

\subsection{Program Pengukuran}

Program berikut menjalankan siklus yang sama dengan program sebelumnya,
tetapi tanpa jendela dan untuk jangka waktu yang tetap. Setiap batas laju
diuji bergantian, lalu hasilnya ditampilkan sebagai tabel.

\kodepython{python/ch04_ukur_laju.py}{Program pengukuran pengaruh batas laju bingkai}{lst:ch04-ukur-laju}

Bagian yang menarik pada program tersebut adalah cara mengukur pemakaian
prosesor. Pemakaian tidak diukur dengan alat luar, melainkan dari dua
penghitung waktu yang disediakan Python.

\begin{itemize}
  \item \code{time.perf_counter()} mengukur \textbf{waktu berjalan}, yaitu lama
        waktu yang dirasakan oleh jam dinding.
  \item \code{time.process_time()} mengukur \textbf{waktu prosesor}, yaitu
        lama prosesor benar-benar bekerja untuk program ini.
\end{itemize}

Perbandingan keduanya menghasilkan pemakaian prosesor. Apabila program
memakai 0,06 detik waktu prosesor selama 2 detik waktu berjalan, pemakaiannya
sekitar tiga persen. Apabila keduanya hampir sama, program memakai satu inti
prosesor secara penuh.

\subsection{Hasil Pengukuran}

Tabel~\ref{tab:hasil-pengukuran-laju} memuat hasil pengukuran pada satu
komputer contoh, dengan lama pengujian dua detik pada setiap konfigurasi.

\begin{table}[htbp]
\centering
\caption{Hasil pengukuran pengaruh batas laju bingkai terhadap pemakaian prosesor}
\label{tab:hasil-pengukuran-laju}
{\setlength{\tabcolsep}{4pt}\footnotesize%
\begin{tabularx}{\textwidth}{@{}C{1.6cm} C{1.8cm} C{2.2cm} C{2.6cm} Y@{}}
\toprule
\headrow \textbf{Batas laju} & \textbf{Bingkai} & \textbf{Laju tercapai} & \textbf{Waktu prosesor} & \textbf{Pemakaian prosesor} \\
\midrule
30 & 60 & 29,7 FPS & 0,02 detik & 0,8\% \\
60 & 120 & 59,9 FPS & 0,06 detik & 3,1\% \\
120 & 233 & 116,0 FPS & 0,09 detik & 4,7\% \\
tanpa & 37.884 & 18.941 FPS & 1,94 detik & 96,9\% \\
\bottomrule
\end{tabularx}}
\tabsource{Pengukuran penulis pada lingkungan tanpa jendela, lama pengujian
2 detik setiap konfigurasi. Angka pada komputer Anda akan berbeda; yang
penting adalah hubungan antarbarisnya.}
\end{table}

Tiga hal dapat dibaca dari tabel tersebut.

\begin{enumerate}
  \item \textbf{Batas laju dipatuhi dengan tepat.} Ketika batasnya 30, laju
        yang tercapai 29,7; ketika batasnya 60, lajunya 59,9. Siklus tidak
        berjalan lebih cepat daripada yang diminta.
  \item \textbf{Pemakaian prosesor tetap kecil pada semua batas yang wajar.}
        Pada 30 sampai 120 bingkai per detik, pemakaiannya berada di bawah
        lima persen. Bahkan pada 120 bingkai per detik, prosesor masih
        menganggur hampir seluruh waktunya.
  \item \textbf{Tanpa pembatasan, prosesor dipakai penuh hampir sepanjang
        waktu.} Perhatikan baris terakhir: 18.941 bingkai per detik berarti
        sekitar enam belas kali lebih banyak daripada yang dapat dilihat mata,
        dan harganya adalah seluruh inti prosesor.
\end{enumerate}

Kesimpulannya sederhana: batas laju bukan pembatasan mutu, melainkan
penghematan yang tidak mengurangi apa pun yang dapat dirasakan pemain.
Permainan yang dibatasi 60 bingkai per detik terlihat sama dengan permainan
yang berjalan ribuan bingkai per detik, tetapi jauh lebih ringan pada
perangkat.

\begin{catatabox}
Angka pada tabel tersebut diukur pada lingkungan tanpa jendela, sehingga
angkanya berbeda dari pengukuran pada jendela sungguhan. Pada jendela
sungguhan, laju tanpa batas biasanya tidak setinggi itu karena menunggu
layar selesai menampilkan gambar, dan pemakaian prosesor juga lebih besar
pada semua baris karena penggambaran sesungguhnya lebih berat daripada
penggambaran di memori. Yang tetap berlaku adalah hubungannya: tanpa batas
laju, prosesor bekerja sekeras mungkin tanpa manfaat yang terlihat.
\end{catatabox}

\subsection{Pilihan Batas Laju}

Dengan bekal pengukuran tersebut, batas laju dapat dipilih secara sadar.

\begin{itemize}
  \item \textbf{60 bingkai per detik} adalah pilihan bawaan yang baik untuk
        sebagian besar permainan sederhana: cukup mulus, dan ringan.
  \item \textbf{30 bingkai per detik} dipakai apabila program mulai berat atau
        perangkatnya sederhana. Gerakannya masih dapat diterima untuk
        permainan yang tidak menuntut ketepatan tinggi.
  \item \textbf{120 bingkai per detik ke atas} hanya bermanfaat apabila
        perangkatnya mendukung dan permainannya menuntut tanggapan yang sangat
        cepat. Untuk keperluan pembelajaran, pilihan ini belum diperlukan.
\end{itemize}

Menariknya, keadaan sebenarnya sering kali berbeda dari dugaan. Batas yang
tinggi tidak dengan sendirinya menghasilkan permainan yang lebih baik,
terutama apabila pekerjaan satu bingkai ternyata memerlukan waktu lebih lama
daripada jatahnya. Contoh pada Bagian~\ref{sec:contoh-bab-4} memperlihatkan
hitungannya.
